\section{Protocole Expérimental}
\subsection{Benchmarks et configuration système}

Afin de comparer les différentes implémentations, nous avant fait 4 benchmarks différents:

Benchmark \textbf{medium}:
\begin{itemize}
	\item 10 primitives
	\item 2 sources de lumières
	\item une bounding box\\
\end{itemize}

Benchmark \textbf{mickey}:
\begin{itemize}
	\item 10 primitives
	\item plusieurs primitives imbriquées
	\item à "ciel ouvert"\\
\end{itemize}

Benchmark \textbf{big}:
\begin{itemize}
	\item 100 primitives
	\item une bounding box\\
\end{itemize}

Benchmark \textbf{huge}:
\begin{itemize}
	\item 1000 primitives
	\item une bounding box\\
\end{itemize}



Les expériences ont été réalisées sur la configuration suivante :
\begin{itemize}
    \item[$\bullet$] CPU : Apple Silicon M4, 4 c\oe urs performance   et 6 c\oe urs efficients
    \item[$\bullet$] Compilateur : \texttt{gcc-15}, flags \texttt{-O3 -march=x86\_64v3}
    \item[$\bullet$] OS : macOS Tahoe 26.4.1
\end{itemize}


\subsection{Mesure du runtime}

Pour ce premier type de protocole, nous allons nous concentrer sur le 
runtime. Il s'agira de mesurer le temps d'exécution de notre programme. 
Nous avons choisi de mesurer exclusivement la phase de rendu, excluant 
ainsi la phase d'initialisation, celle de l'écriture de l'image et du 
calcul du runtime lui-même.\\

L'objectif sera de comparer le runtime en faisant varier différents 
paramètres, qui seront explicités pour chaque série de mesures. La scène 
de référence utilisée est fixée à une résolution de $800 \times 600$ 
pixels, avec $[\ldots]$ primitives, $[\ldots]$ samples et $[\ldots]$ 
rebonds maximum.\\

Afin de limiter les incertitudes systématiques, nous avons utilisé la 
fonction \texttt{clock\_gettime(CLOCK\_MONOTONIC)} de \texttt{<time.h>}, 
qui mesure le temps réel écoulé avec une résolution de $10^{-9}$ secondes, 
réduisant ainsi drastiquement l'incertitude de l'outil de mesure. Chaque 
mesure est répétée 10 fois ; nous retenons la moyenne des répétitions, 
après suppression des valeurs aberrantes éventuelles.\\

Critères de comparaison : on considèrera qu'une optimisation est 
bénéfique sur le runtime si le gain mesuré est supérieur à $5\%$.

\subsection{Mesure de la convergence}

Pour ce second type de protocole, nous allons mesurer la convergence 
visuelle du rendu en fonction du nombre de samples. L'objectif est de 
quantifier l'écart entre une image rendue avec $N$ samples et une image 
de référence, considérée comme vraie.\\

\subsubsection{Image de référence}

L'image de référence est obtenue en rendant la scène avec un nombre de 
samples $N_{ref}$ suffisamment grand pour que la variance résiduelle soit négligeable devant les écarts que nous souhaitons mesurer. Nous fixons 
$N_{ref} = 10000$ samples pour chaque benchmarks.\\

\subsubsection{Métrique de convergence}

Pour comparer une image rendue $I_N$ à l'image de référence $I_{ref}$, 
nous utilisons l'erreur quadratique moyenne (RMSE, \textit{Root Mean 
Squared Error}) :

$$\text{RMSE}(N) = \sqrt{\frac{1}{3 \cdot W \cdot H} \sum_{i=1}^{W} 
\sum_{j=1}^{H} \sum_{c \in \{R,G,B\}} 
\left(I_N(i,j,c) - I_{ref}(i,j,c)\right)^2}$$

où $W$ et $H$ sont respectivement la largeur et la hauteur de l'image 
en pixels. Le facteur $3$ provient de la sommation sur les trois canaux 
de couleur. Les valeurs de pixels sont normalisées dans $[0, 1]$.\\

Le RMSE est préféré au MSE car il s'exprime dans la même unité que les 
valeurs de pixels, facilitant l'interprétation. On s'attend théoriquement 
à ce que le RMSE décroisse en $O\left(\frac{1}{\sqrt{N}}\right)$, 
conformément à la convergence de l'estimateur Monte Carlo.\\

\subsubsection{Déroulement des mesures}

Pour chaque configuration testée (naïve, BVH, clustering, SIMD), on rend la 
scène pour $N \in \{10, 50, 100, 200, 500, 1000, 2000\}$ samples. Le 
script Python calcule le RMSE entre chaque image et $I_{ref}$, puis trace 
$\text{RMSE}(N)$ en échelle log-log afin de visualiser la pente de 
convergence. Une pente de $-\frac{1}{2}$ en log-log confirme la 
convergence théorique en $O\left(\frac{1}{\sqrt{N}}\right)$.